\documentclass[11pt]{article}

\usepackage[margin=1in]{geometry}
\usepackage{amsmath}
\usepackage{amssymb}
\usepackage{amsthm}
\usepackage{mathtools}
\usepackage{xcolor}
\usepackage[colorlinks=true,linkcolor=blue,urlcolor=blue]{hyperref}
\usepackage[final]{microtype}

\title{heterosplat: Kernel Mathematics}
\author{Source: \texttt{repos/heterosplat/Documents/LaTeX/KernelMathematics.tex}}
\date{\today}

\newcommand{\SO}{\mathrm{SO}}
\newcommand{\SU}{\mathrm{SU}}
\newcommand{\SymPos}{\mathrm{Sym}^{+}}
\newcommand{\diag}{\operatorname{diag}}
\newcommand{\tr}{\operatorname{tr}}
\newcommand{\R}{\mathbb{R}}
\newcommand{\HQ}{\mathbb{H}}
\newcommand{\PSdef}{\succ 0}

\begin{document}
\maketitle

\section{Purpose and audience}

This document spells out the mathematics of each CUDA kernel that
heterosplat ships. Each section names the source files it documents
and gives the equations actually computed there. The intended reader
is comfortable with elementary differential geometry (smooth maps
between manifolds, Jacobians of matrix-valued functions, the Lie
group $\SO(3)$) and abstract algebra (group actions, quotient
spaces). The document is meant to be read alongside the source so
the reader can verify the implementation against the math.

\section{Conventions}

\paragraph{Indexing.}
Vectors are 1-indexed in math, $0$-indexed in code; we write the
correspondence explicitly when relevant. All matrices are stored as
contiguous arrays of \texttt{float}; full $3 \times 3$ blocks are
\emph{row-major} on disk, while the GLM library's in-register type
\texttt{glm::mat3} is \emph{column-major}. Kernels that emit a full
covariance therefore perform a transpose-on-write
(\texttt{covars[i*3+j] = covar[j][i]}).

\paragraph{Symmetric storage.}
A symmetric $3\times 3$ matrix may be stored as either:
\begin{itemize}
  \item the full $9$-float row-major layout; or
  \item a $6$-float upper-triangle (``triu'') layout in the order
  $\bigl(\Sigma_{11},\,\Sigma_{12},\,\Sigma_{13},\,\Sigma_{22},\,\Sigma_{23},\,\Sigma_{33}\bigr)$.
\end{itemize}
The kernel's \texttt{triu} boolean selects between them.

\paragraph{Quaternion order.}
A quaternion is written $q = (w, x, y, z) \in \HQ$ with the real part
\emph{first}; this matches gsplat's storage. Unit quaternions
$\HQ_{1} \cong \SU(2)$ double-cover $\SO(3)$ via $q \mapsto R(q)$;
both $q$ and $-q$ map to the same rotation.

\paragraph{Vector-Jacobian product (VJP).}
For a smooth map $f : \R^{n} \to \R^{m}$ and a covector
$v \in (\R^{m})^{*}$, the VJP is $\mathrm{VJP}_{f}(x; v) = v^{T}\,
\partial f / \partial x \in (\R^{n})^{*}$. In machine-learning idiom,
if $L$ is a downstream scalar loss and $v = \partial L/\partial f$,
then $\mathrm{VJP}_{f}(x; v) = \partial L/\partial x$. Every backward
kernel in heterosplat is a VJP.

\section{Quaternion + scale \texorpdfstring{$\to$}{to} covariance / precision}
\label{sec:qsc}

\paragraph{Source files.}
\begin{itemize}
  \item Vendored kernel templates (extracted from upstream gsplat):\\
  \texttt{Source/Kernels/Thirdparty/Gsplat/QuatScaleToCovarKernels.cuh}
  \item Vendored device helpers used by the kernels:\\
  \texttt{Source/Kernels/Thirdparty/Gsplat/Utils.cuh} ---
  specifically \texttt{quat\_to\_rotmat}, \texttt{quat\_to\_rotmat\_vjp},
  \texttt{quat\_scale\_to\_covar\_preci}, \texttt{quat\_scale\_to\_covar\_vjp},
  \texttt{quat\_scale\_to\_preci\_vjp}.
  \item Our torch-free launcher:\\
  \texttt{Source/Kernels/Heterosplat/QuatScaleToCovar.\{h,cu\}}.
  \item Smoke tests:
  \texttt{Source/UnitTests/Kernels/Heterosplat/QuatScaleToCovar\_tests.cu}.
\end{itemize}

\subsection{The forward map}

A 3D Gaussian splat is parameterized by a quaternion
$q \in \HQ \setminus \{0\}$ and a positive scale vector
$s = (s_{1}, s_{2}, s_{3}) \in \R^{3}_{>0}$. From these, build
\begin{equation}
  \hat q \;=\; \frac{q}{\lVert q \rVert}, \qquad
  R \;=\; R(\hat q) \in \SO(3), \qquad
  S \;=\; \diag(s_{1},\,s_{2},\,s_{3}) \in \SymPos(3).
\end{equation}
The covariance and its inverse (the \emph{precision}) are
\begin{align}
  \Sigma &\;=\; R\, S^{2}\, R^{T} \;\in\; \SymPos(3),
  \label{eq:cov}\\
  \Pi &\;=\; \Sigma^{-1} \;=\; R\, S^{-2}\, R^{T} \;\in\; \SymPos(3).
  \label{eq:prec}
\end{align}
Equivalently, defining $M = R S$,
\begin{equation}
  \Sigma \;=\; M M^{T}, \qquad
  M' \;=\; R S^{-1} \,\Longrightarrow\, \Pi \;=\; M' (M')^{T}.
  \label{eq:m-form}
\end{equation}
The kernel uses the $M$-form (\eqref{eq:m-form}) because it gives a
single matrix product per Gaussian regardless of which output
($\Sigma$, $\Pi$, or both) is requested.

\paragraph{Quaternion to rotation.}
For $q = (w, x, y, z)$ with $\lVert q \rVert^{2} = w^{2} + x^{2} + y^{2} + z^{2}$
and $\hat q = q / \lVert q \rVert = (\hat w, \hat x, \hat y, \hat z)$,
\begin{equation}
  R(q) \;=\;
  \begin{pmatrix}
    1 - 2(\hat y^{2} + \hat z^{2}) &
      2(\hat x \hat y - \hat w \hat z) &
      2(\hat x \hat z + \hat w \hat y) \\[2pt]
    2(\hat x \hat y + \hat w \hat z) &
      1 - 2(\hat x^{2} + \hat z^{2}) &
      2(\hat y \hat z - \hat w \hat x) \\[2pt]
    2(\hat x \hat z - \hat w \hat y) &
      2(\hat y \hat z + \hat w \hat x) &
      1 - 2(\hat x^{2} + \hat y^{2})
  \end{pmatrix}.
  \label{eq:quat-to-rotmat}
\end{equation}
This is implemented in \texttt{quat\_to\_rotmat}. Inline
normalisation ensures $R \in \SO(3)$ for any $q \neq 0$, so the
optimizer is free to update $q$ off the unit sphere; the geometry is
$\mathbb{H}_{1}$-equivariant.

\paragraph{Geometric interpretation.}
The map
\begin{equation}
  \Phi \,\colon\, \SO(3) \times \R^{3}_{>0} \longrightarrow \SymPos(3),
  \qquad (R, s) \longmapsto R\, \diag(s)^{2}\, R^{T}
\end{equation}
is a smooth surjection. On the open dense set of SPD matrices with
distinct eigenvalues the map is a submersion; the fiber over a generic
$\Sigma$ is the orbit of the \emph{signed-permutation group}
$B_{3} \le \mathrm{O}(3)$ acting by
$(R, s) \mapsto (R P,\, P^{T} s)$ where $P$ is a $3 \times 3$ signed
permutation matrix and $P^{T} s$ permutes the (positive) entries.
$B_{3}$ has order $48$. Composed with the $\HQ_{1} \to \SO(3)$ double
cover, the $(q, s)$ parameterization is $96$-fold redundant. The
splat loss is invariant under this redundancy, so optimisation does
not need to break it.

\paragraph{Why this parameterization.}
A direct optimisation of $\Sigma$ over $\SymPos(3)$ requires a
Riemannian retraction (or a Cholesky lift with positivity
constraints). The $(q, s)$ form sidesteps this:
$\HQ \setminus \{0\}$ is a $4$-dim open subset of $\R^{4}$ and
$\R^{3}_{>0}$ is a $3$-dim open submanifold of $\R^{3}$, both
log-trivial; ordinary Adam works on each. Eigenvalues of $\Sigma$
are exactly $s_{i}^{2}$, so the principal-axis interpretation is
preserved at all times.

\subsection{Storage layout (forward output)}

For a batch of $N$ Gaussians with output pointer $\Sigma_{\text{out}}$:
\begin{itemize}
  \item Full layout (\texttt{triu = false}): $\Sigma_{\text{out}}[N, 3, 3]$
  is row-major; entry $(i, j)$ of Gaussian $n$ lives at byte index
  $(9n + 3i + j) \cdot \texttt{sizeof(float)}$. Because GLM's
  \texttt{mat3} is column-major, the kernel writes
  \texttt{covars[i*3 + j] = covar[j][i]}.
  \item Triu layout (\texttt{triu = true}): the six independent
  entries are written in row-major upper-triangle order, i.e.\
  $(\Sigma_{11}, \Sigma_{12}, \Sigma_{13}, \Sigma_{22}, \Sigma_{23}, \Sigma_{33})$.
\end{itemize}

\subsection{Launcher API \texorpdfstring{$\leftrightarrow$}{<->} math correspondence (forward)}
\label{sec:qsc-fwd-api}

The C++ launcher
\texttt{launch\_quat\_scale\_to\_covar\_preci\_forward} (in
\texttt{Source/Kernels/Heterosplat/QuatScaleToCovar.\{h,cu\}}) is a thin
host wrapper that fires the templated \texttt{\_\_global\_\_} kernel
once per Gaussian. Each parameter maps to a math object as follows.

\begin{center}
\begin{tabular}{l l p{8.0cm}}
\hline
\textbf{C++ parameter} & \textbf{Math symbol} & \textbf{Meaning} \\
\hline
\texttt{N} & $N$ &
Batch size; one CUDA thread per Gaussian, blocks of $256$. \\
\texttt{quats} & $\{q_n\}_{n=1}^{N}$ &
Device buffer $[N, 4]$, row $n$ is $(w_n, x_n, y_n, z_n)$ (real part first).
Need not be unit length: kernel normalises in-line. \\
\texttt{scales} & $\{s_n\}_{n=1}^{N}$ &
Device buffer $[N, 3]$, row $n$ is $(s_1^{(n)}, s_2^{(n)}, s_3^{(n)})$
with $s_k^{(n)} > 0$. \\
\texttt{triu} & --- &
Output layout flag: \texttt{true} = $[N,6]$ upper-triangle order,
\texttt{false} = $[N,9]$ row-major full $3 \times 3$. \\
\texttt{covars} & $\{\Sigma_n\}_{n=1}^{N}$ &
Device output for $\Sigma_n = R_n S_n^{2} R_n^T$. \texttt{nullptr}
skips the covariance branch entirely. \\
\texttt{precis} & $\{\Pi_n\}_{n=1}^{N}$ &
Device output for $\Pi_n = R_n S_n^{-2} R_n^T$. \texttt{nullptr}
skips the precision branch. \\
\texttt{stream} & --- &
CUDA stream for asynchronous overlap; \texttt{nullptr} = default stream. \\
\hline
\end{tabular}
\end{center}

\noindent
The launch is a no-op when $N = 0$ or both \texttt{covars} and
\texttt{precis} are null. The kernel writes (does not accumulate);
prior contents of the output buffers are overwritten.

\subsection{The backward map (VJP)}
\label{sec:qsc-bwd}

Suppose we receive an upstream covariance gradient
\begin{equation}
  G \;=\; \frac{\partial L}{\partial \Sigma} \;\in\; \R^{3 \times 3},
\end{equation}
and we want to push it back to $(q, s)$. Two key matrix-calculus
identities do the work.

\paragraph{Identity 1: \texorpdfstring{$\Sigma = M M^{T}$}{Sigma = M M\^{}T}.}
For $\Sigma = M M^{T}$ with $M \in \R^{3 \times 3}$,
\begin{equation}
  \frac{\partial L}{\partial M} \;=\; \bigl(G + G^{T}\bigr)\, M.
  \label{eq:vjp-mmt}
\end{equation}
Derivation: $\partial \Sigma_{ij} / \partial M_{kl} =
\delta_{ik} M_{jl} + \delta_{jk} M_{il}$, so
$\partial L / \partial M_{kl} = \sum_{ij} G_{ij}
(\delta_{ik} M_{jl} + \delta_{jk} M_{il})
= (GM)_{kl} + (G^{T} M)_{kl}$.

\paragraph{Identity 2: \texorpdfstring{$M = R S$}{M = R S} with \texorpdfstring{$S$}{S} diagonal.}
\begin{align}
  \frac{\partial L}{\partial R} &\;=\;
  \frac{\partial L}{\partial M}\, S, \label{eq:vjp-R}\\
  \left(\frac{\partial L}{\partial s}\right)_{k} &\;=\;
  \sum_{i = 1}^{3} R_{ik}\, \biggl(\frac{\partial L}{\partial M}\biggr)_{ik}
  \quad \text{for } k \in \{1, 2, 3\}. \label{eq:vjp-s}
\end{align}
\eqref{eq:vjp-R} follows from $M = R S$ and $S$ being on the right;
\eqref{eq:vjp-s} comes from $M_{ij} = \sum_{k} R_{ik} S_{kj}
= R_{ij} s_{j}$ (since $S$ is diagonal), so
$\partial M_{ij} / \partial s_{k} = R_{ik}\, \delta_{jk}$.

\paragraph{Putting it together.}
Combining \eqref{eq:vjp-mmt}--\eqref{eq:vjp-s} gives the closed form
implemented in \texttt{quat\_scale\_to\_covar\_vjp}:
\begin{align}
  \widetilde G &\;=\; G + G^{T}, \\
  \frac{\partial L}{\partial M} &\;=\; \widetilde G \, M, \\
  \frac{\partial L}{\partial R} &\;=\; \widetilde G \, M \, S \;=\; \widetilde G \, R \, S^{2}, \\
  \biggl(\frac{\partial L}{\partial s}\biggr)_{k} &\;=\; \sum_{i} R_{ik} \,(\widetilde G M)_{ik}.
\end{align}
Finally $\partial L / \partial q$ comes from $\partial L / \partial R$
via the quaternion-to-rotation Jacobian, implemented in
\texttt{quat\_to\_rotmat\_vjp}. Because $R$ uses the normalised
$\hat q = q / \lVert q \rVert$, that helper additionally projects
the gradient onto $T_{q}\HQ_{1}$ (the tangent space at $q$) so
non-unit updates do not spuriously enlarge $q$ along its own
direction.

\paragraph{Triu input gradients (the $\tfrac{1}{2}$ factor).}
When the upstream gradient is supplied as a $6$-vector $g$ in the
triu layout, expanding to a symmetric $3 \times 3$ matrix
double-counts the off-diagonals (each $g_{12}$ contributes once to
$\Sigma_{12}$ \emph{and} once to $\Sigma_{21}$, but only one of
those is stored). The kernel halves the off-diagonals:
\begin{equation}
  G \;=\; \begin{pmatrix}
    g_{1} & \tfrac{1}{2} g_{2} & \tfrac{1}{2} g_{3} \\
    \tfrac{1}{2} g_{2} & g_{4} & \tfrac{1}{2} g_{5} \\
    \tfrac{1}{2} g_{3} & \tfrac{1}{2} g_{5} & g_{6}
  \end{pmatrix}.
  \label{eq:triu-half}
\end{equation}
This is the literal $\texttt{0.5f}$ multiplier visible in
\texttt{quat\_scale\_to\_covar\_preci\_bwd\_kernel} when
\texttt{triu == true}.

\subsection{Launcher API \texorpdfstring{$\leftrightarrow$}{<->} math correspondence (backward)}
\label{sec:qsc-bwd-api}

The companion launcher
\texttt{launch\_quat\_scale\_to\_covar\_preci\_backward} carries the
gradient through the operation in §\ref{sec:qsc-bwd}.

\paragraph{The \texttt{v\_} prefix.}
Every gradient parameter is named \texttt{v\_X}. The convention is
inherited from gsplat / common automatic-differentiation idiom: read
\texttt{v\_X} as ``the value of the gradient flowing through $X$,''
i.e.\ $\partial L / \partial X$. The letter $v$ is short for
\emph{vector-Jacobian product}; in the language of differential
geometry, $v_X$ is the cotangent (covector) dual to a tangent
direction $\mathrm{d}X$, pulled back from the loss.

\paragraph{Why the forward inputs are passed in again.}
The PyTorch path stashes $q$ and $s$ in autograd's saved-tensor cache
during the forward call, so they are silently available for backward.
heterosplat has no such cache: the caller provides them again, and the
kernel locally recomputes $R(q)$ and $M = R S$ to apply the chain rule.
This is an explicit memory-versus-compute trade. Because the
$\Theta(N)$ recompute is dominated by the rasterizer's
$\Theta(N \cdot \text{tiles})$ work later in the pipeline, the wall-time
cost is invisible in practice.

\begin{center}
\begin{tabular}{l l p{8.0cm}}
\hline
\textbf{C++ parameter} & \textbf{Math object} & \textbf{Meaning} \\
\hline
\texttt{N} & $N$ &
Batch size; must equal the value passed to the forward launcher. \\
\texttt{quats}, \texttt{scales} & $\{q_n\}, \{s_n\}$ &
Same forward inputs, used to re-derive $R_n$ and $M_n$ in-kernel. \\
\texttt{triu} & --- &
Layout of the input gradient buffers; must match what the forward call
emitted. When \texttt{true}, the kernel applies the $\tfrac{1}{2}$
factor on off-diagonals from~\eqref{eq:triu-half}. \\
\texttt{v\_covars} & $\{ \partial L / \partial \Sigma_n \}$ &
\emph{Input} gradient. \texttt{nullptr} signals ``$\Sigma$ did not feed
the loss''; the covariance-side chain rule is skipped. \\
\texttt{v\_precis} & $\{ \partial L / \partial \Pi_n \}$ &
\emph{Input} gradient for the precision branch; same \texttt{nullptr}
semantics as \texttt{v\_covars}. \\
\texttt{v\_quats} & $\{ \partial L / \partial q_n \}$ &
\emph{Output} gradient, $[N, 4]$. The kernel \emph{writes} (does not
accumulate). Caller is responsible for any mixing with prior gradients. \\
\texttt{v\_scales} & $\{ \partial L / \partial s_n \}$ &
\emph{Output} gradient, $[N, 3]$; same write-not-accumulate semantics. \\
\texttt{stream} & --- &
CUDA stream; \texttt{nullptr} = default. \\
\hline
\end{tabular}
\end{center}

\noindent
\emph{Caveat:} when both \texttt{v\_covars} and \texttt{v\_precis} are
\texttt{nullptr} the launch is skipped \emph{and the output buffers are
left untouched}. If ``no gradient'' should mean a literal zero
(rather than the previous training step's stale memory), the caller
must zero \texttt{v\_quats} and \texttt{v\_scales} themselves.

\subsection{The precision branch}

For $\Pi = R S^{-2} R^{T}$, the same derivation goes through with
$M' = R S^{-1}$ in place of $M$, plus a chain-rule factor for the
inversion of $s$:
\begin{equation}
  \frac{\partial s_{k}^{-1}}{\partial s_{k}} \;=\; -\frac{1}{s_{k}^{2}}.
\end{equation}
The kernel asserts $s_{k} \neq 0$ before forming $1 / s_{k}$. In the
production pipeline, frustum culling and minimum-radius filters drop
degenerate Gaussians upstream, so this assertion is a defensive
check and never fires in practice.

\subsection{Numerical worth-noting}

\begin{itemize}
  \item The forward kernel uses \texttt{rsqrt} for quaternion
  normalisation rather than $1 / \sqrt{\cdot}$; on modern GPUs this
  is a single intrinsic with relative error $\le 2^{-23}$.
  \item Inputs and outputs are all \texttt{float32}; double-precision
  accumulation is not used. For $N \lesssim 10^{6}$ Gaussians (the
  typical splat scene) this is well within tolerance.
  \item The forward kernel is parallel over $N$ with no shared
  memory; one Gaussian per thread, $256$ threads per block. The
  backward is identical in structure.
\end{itemize}

\subsection{Test correspondence}

The smoke tests in \texttt{QuatScaleToCovar\_tests.cu} cover three
identities from \eqref{eq:cov}:
\begin{enumerate}
  \item $q = (1, 0, 0, 0)$, $s = (1, 1, 1)$
  $\Longrightarrow$
  $\Sigma = R \, S^{2} \, R^{T} = I_{3} \cdot I_{3} \cdot I_{3}^{T} = I_{3}$.
  \item $q = (1, 0, 0, 0)$, $s = (2, 3, 5)$
  $\Longrightarrow$
  $\Sigma = \diag(4, 9, 25)$.
  \item Triu layout from the same input matches the diagonal entries
  in positions $1, 4, 6$ of the $6$-vector and zeros elsewhere
  (verifying the storage convention from \eqref{eq:triu-half} on the
  forward side).
\end{enumerate}

\section{Spherical harmonics (view-dependent colour)}
\label{sec:sh}

\paragraph{Source files.}
\begin{itemize}
  \item Vendored kernel templates and \texttt{\_\_device\_\_} helpers:\\
  \texttt{Source/Kernels/Thirdparty/Gsplat/SphericalHarmonicsKernels.cuh}
  (extracted from upstream gsplat's \texttt{SphericalHarmonicsCUDA.cu};
  the only patch is replacing ATen's \texttt{gpuAtomicAdd} with the
  CUDA built-in \texttt{atomicAdd}).
  \item Our torch-free launchers:
  \texttt{Source/Kernels/Heterosplat/SphericalHarmonics.\{h,cu\}}.
  \item Smoke tests:
  \texttt{Source/UnitTests/Kernels/Heterosplat/SphericalHarmonics\_tests.cu}.
\end{itemize}

\subsection{The forward map}

Each Gaussian carries a per-channel set of real spherical-harmonic
(SH) coefficients
$\{c_{l,m}^{(\text{ch})}\}_{l = 0}^{L},\, m = -l, \dots, l$ for
$\text{ch} \in \{R, G, B\}$. Given a unit view direction
$\hat d \in S^2$, the kernel emits a colour
\begin{equation}
  \text{colour}^{(\text{ch})}(\hat d)
  \;=\;
  \sum_{l=0}^{L} \sum_{m=-l}^{l}
    c_{l,m}^{(\text{ch})}\, Y_{l,m}(\hat d),
\end{equation}
where $Y_{l,m} : S^2 \to \R$ are the real-valued spherical harmonics.
The kernel parameter \texttt{degrees\_to\_use} fixes $L$ at runtime so
coarse-to-fine training (start at $L = 0$, lift one degree at a time)
costs no re-allocation.

\paragraph{Background.}
The $Y_{l,m}$ are an orthonormal basis of $L^{2}(S^{2})$ with respect
to the surface-area measure: they are eigenfunctions of the
Laplace--Beltrami operator $\Delta_{S^{2}}$ with eigenvalue $-l(l+1)$.
The complex spherical harmonics $Y_{l}^{m}$ admit several "real-basis"
conventions; gsplat (and therefore heterosplat) uses Sloan's
fast-evaluation form (JCGT 2013), which embeds normalisation
constants into the recursion and absorbs Condon--Shortley sign
factors. The first two degrees, in the order the kernel walks
coefficient indices $0, 1, 2, 3$:
\begin{align}
  Y_{0,0}(\hat d) &= \tfrac{1}{2\sqrt{\pi}} \;\approx\; 0.28209479\dots, \\
  Y_{1,-1}(\hat d) &= -\sqrt{\tfrac{3}{4\pi}}\, y \;\approx\; -0.48860251\dots\, y, \\
  Y_{1,0}(\hat d)  &= \phantom{-}\sqrt{\tfrac{3}{4\pi}}\, z \;\approx\; \phantom{-}0.48860251\dots\, z, \\
  Y_{1,1}(\hat d)  &= -\sqrt{\tfrac{3}{4\pi}}\, x \;\approx\; -0.48860251\dots\, x,
\end{align}
which is the literal three-term degree-1 expansion in
\texttt{sh\_coeffs\_to\_color\_fast}. Higher-degree terms follow the
\emph{associated Legendre} recursion in normalised form; the
embedded floats (\texttt{0.94617469}\dots etc.) are precomputed
products of $\sqrt{(2l+1)/(4\pi)}$ and Condon--Shortley factors.

\paragraph{Direction normalisation.}
The kernel does \emph{not} require its callers to pre-normalise
\texttt{dirs}; it computes
$\text{inorm} = 1 / \sqrt{\hat d \cdot \hat d}$ inline and uses
$\hat d \cdot \text{inorm}$ throughout the recursion. This is the
correct form because $Y_{l,m}$ is defined on the unit sphere; passing
in a non-unit vector is equivalent to evaluating at its radial
projection.

\paragraph{Per-Gaussian masking.}
The optional \texttt{masks} buffer is a \texttt{[N]} array of
\texttt{bool}; threads with \texttt{masks[n] == false} short-circuit
\emph{after} the SH ID dispatch, leaving the corresponding output
slot untouched. Useful for selectively training subsets of the
splat without re-allocating coefficient buffers.

\subsection{Storage layout}

\begin{itemize}
  \item \texttt{dirs} -- \texttt{[N, 3]} row-major, contiguous floats.
  Internally cast to \texttt{vec3*}; the layout is identical.
  \item \texttt{coeffs} -- \texttt{[N, K, 3]} row-major. \texttt{K} must
  be at least $(L+1)^{2}$; e.g.\ degree $0 \to K \ge 1$, degree
  $1 \to K \ge 4$, ..., degree $4 \to K \ge 25$. The kernel uses
  \texttt{K} as the per-Gaussian stride only, so \texttt{K} can exceed
  $(L+1)^{2}$ if the caller wants to keep a degree-4 buffer
  pre-allocated while training at lower $L$.
  \item \texttt{colors} -- \texttt{[N, 3]} row-major, RGB.
\end{itemize}

\subsection{Launcher API \texorpdfstring{$\leftrightarrow$}{<->} math correspondence (forward)}

\begin{center}
\begin{tabular}{l l p{8.0cm}}
\hline
\textbf{C++ parameter} & \textbf{Math symbol} & \textbf{Meaning} \\
\hline
\texttt{N} & $N$ & Number of Gaussians; one CUDA thread per
\emph{(Gaussian, channel)} pair (so total threads = $3N$). \\
\texttt{K} & $K$ & Per-Gaussian coefficient stride; $K \ge (L+1)^{2}$. \\
\texttt{degrees\_to\_use} & $L$ & Maximum SH degree to evaluate (0--4
supported). \\
\texttt{dirs} & $\{\hat d_n\}$ & View directions, normalised in-kernel. \\
\texttt{coeffs} & $\{c_{l,m,n}^{(\text{ch})}\}$ & SH coefficients. \\
\texttt{masks} & --- & Optional per-Gaussian skip flag;
\texttt{nullptr} = process all. \\
\texttt{colors} & $\{\text{colour}_n\}$ & Output, $[N, 3]$. \\
\texttt{stream} & --- & CUDA stream. \\
\hline
\end{tabular}
\end{center}

\subsection{The backward map (VJPs)}

Backward propagates $v_{\text{colors}} = \partial L / \partial \text{colours}$
to two parameters: the SH coefficients and (optionally) the view
direction.

\paragraph{Coefficient gradient (linear).}
Because the forward is linear in coefficients, the VJP is closed
form:
\begin{equation}
  \frac{\partial L}{\partial c_{l,m,n}^{(\text{ch})}}
  \;=\;
  Y_{l,m}(\hat d_n) \cdot v_{\text{colors}}^{(n,\, \text{ch})}.
\end{equation}
The kernel writes each \texttt{v\_coeffs[n, l m, ch]} slot \emph{once}
(single-writer, no atomics).

\paragraph{Direction gradient (chain rule with unit-sphere projection).}
The dependence on $\hat d$ is more interesting because $\hat d \in S^{2}$.
Internally the kernel works with the normalised
$\hat d = d / \lVert d \rVert$ and we want $\partial L / \partial d$
(the gradient w.r.t.\ the \emph{raw}, possibly non-unit, input).

Computing $\partial L / \partial \hat d$ in the rotated frame is
straightforward (each coefficient contributes a polynomial in
$\hat x, \hat y, \hat z$). Pulling that back to $d$ requires the
unit-sphere projection
\begin{equation}
  \frac{\partial L}{\partial d}
  \;=\;
  \frac{1}{\lVert d \rVert}
  \left(I - \hat d\, \hat d^{T}\right)
  \frac{\partial L}{\partial \hat d},
  \label{eq:sphere-projection}
\end{equation}
which is exactly what the line
\texttt{(v\_dir\_n - dot(v\_dir\_n, dir\_n) * dir\_n) * inorm}
computes: $(I - \hat d \hat d^{T})$ removes the radial component
(motion that does not stay on $S^{2}$), and the $1/\lVert d \rVert$
factor converts the on-sphere covector to one on the ambient
$\R^{3}$.

\paragraph{Atomic accumulation.}
Each Gaussian is touched by three threads (one per colour channel).
The coefficient gradient is per-(Gaussian, basis, channel), so each
slot has a single writer. The direction gradient, in contrast, is
per-Gaussian only -- all three channels write into
$\partial L / \partial d_n$. The kernel uses \texttt{atomicAdd} on the
three components to merge the contributions safely. Caller must
\textbf{pre-zero} \texttt{v\_dirs}; the kernel only accumulates.

\subsection{Launcher API \texorpdfstring{$\leftrightarrow$}{<->} math correspondence (backward)}

\begin{center}
\begin{tabular}{l l p{8.0cm}}
\hline
\textbf{C++ parameter} & \textbf{Math object} & \textbf{Meaning} \\
\hline
\texttt{N}, \texttt{K}, \texttt{degrees\_to\_use} & --- &
Same shape parameters as forward. \\
\texttt{dirs}, \texttt{coeffs} & $\hat d, c$ &
Forward inputs; needed for the VJP chain rule. \\
\texttt{masks} & --- & Optional skip flag (same as forward). \\
\texttt{v\_colors} & $\partial L / \partial \text{colours}$ &
\emph{Input} gradient, $[N, 3]$. \\
\texttt{v\_coeffs} & $\partial L / \partial c$ &
\emph{Output} gradient, $[N, K, 3]$. Single-writer; kernel
\emph{writes} (no accumulate). \\
\texttt{v\_dirs} & $\partial L / \partial d$ &
\emph{Output} gradient, $[N, 3]$. \texttt{nullptr} skips the path
entirely; otherwise the kernel \emph{accumulates} via
\texttt{atomicAdd} -- caller must pre-zero. \\
\texttt{stream} & --- & CUDA stream. \\
\hline
\end{tabular}
\end{center}

\subsection{Test correspondence}

The smoke tests in \texttt{SphericalHarmonics\_tests.cu} cover:
\begin{enumerate}
  \item \emph{Degree 0 is DC}: with $L = 0$ the only basis is
  $Y_{0,0} = 1/(2\sqrt{\pi})$, so
  $\text{colours}_n = 0.282\dots \cdot c_{0, n}$.
  \item \emph{Degree 1 single-basis}: with $\hat d = (0, 1, 0)$ and
  only $c_{1,-1}$ non-zero, the kernel reduces to
  $\text{colours} = -0.488\dots \cdot c_{1,-1}$, isolating the
  $Y_{1,-1}$ term.
  \item \emph{Mask short-circuit}: a \texttt{masks[n] = false}
  Gaussian's output slot is left at a sentinel value, confirming the
  early-return path.
  \item \emph{Backward gradcheck (coefficients)}: the analytic
  \texttt{v\_coeffs} from the kernel matches a centered
  finite-difference numerical gradient through the forward to within
  $\text{atol} = 10^{-3}$ + $\text{rtol} = 5 \times 10^{-3}$.
\end{enumerate}

\section{Tile intersection}
\label{sec:intersect-tile}

\paragraph{Source files.}
\begin{itemize}
  \item Vendored kernel template and AccuTile/SNUGBOX helpers
  (extracted from upstream gsplat):\\
  \texttt{Source/Kernels/Thirdparty/Gsplat/IntersectTileKernels.cuh}
  \item Raw-pointer launcher owned by heterosplat:\\
  \texttt{Source/Kernels/Heterosplat/IntersectTile.\{h,cu\}}
  \item Tests and oracle fixtures:\\
  \texttt{Source/UnitTests/Kernels/Heterosplat/IntersectTile\_tests.cu}\\
  \texttt{Source/UnitTests/Kernels/Heterosplat/IntersectTile\_oracle\_tests.cu}\\
  \texttt{Source/UnitTests/Fixtures/IntersectTile/}
\end{itemize}

\subsection{The forward map}

The tile-intersection kernel converts projected 2D Gaussians into a
flat list of tile intersections. For each image $i \in \{0,\dots,I-1\}$
and Gaussian $n \in \{0,\dots,N-1\}$, the input is:
\[
  \mu_{i,n} = (x_{i,n}, y_{i,n}) \in \R^{2}, \qquad
  r_{i,n} = (r^x_{i,n}, r^y_{i,n}) \in \mathbb{Z}^{2}, \qquad
  z_{i,n} \in \R.
\]
In packed mode, the dense pair $(i,n)$ is replaced by a packed row
$p \in \{0,\dots,\mathrm{nnz}-1\}$ and an explicit image id
\texttt{image\_ids[p]}. The output \texttt{flatten\_ids} always stores
the dense flatten index $iN+n$ or the packed row $p$.

The kernel is a two-pass map. First it counts how many tiles each
Gaussian touches:
\[
  c_k = \lvert T_k \rvert,
\]
where $k$ is either the dense flatten index or packed row, and $T_k$
is the set of intersected tile ids. The caller prefix-sums these
counts to obtain
\[
  C_k = \sum_{j=0}^{k} c_j.
\]
The second pass writes the actual intersections into the half-open
slice
\[
  [\, C_{k-1}, C_k \,),
  \qquad C_{-1} = 0.
\]

\paragraph{AABB fallback.}
When \texttt{conics == nullptr} or \texttt{opacities == nullptr}, the
kernel uses the axis-aligned radius rectangle. Let $s$ be
\texttt{tile\_size}; tile coordinates are
\[
  u = x/s, \qquad v = y/s, \qquad
  \rho_x = r^x/s, \qquad \rho_y = r^y/s.
\]
The inclusive/exclusive tile bounds are:
\[
\begin{aligned}
  t^x_{\min} &= \operatorname{clamp}(\lfloor u-\rho_x \rfloor, 0, W), &
  t^x_{\max} &= \operatorname{clamp}(\lceil u+\rho_x \rceil, 0, W),\\
  t^y_{\min} &= \operatorname{clamp}(\lfloor v-\rho_y \rfloor, 0, H), &
  t^y_{\max} &= \operatorname{clamp}(\lceil v+\rho_y \rceil, 0, H),
\end{aligned}
\]
where $W=\texttt{tile\_width}$ and $H=\texttt{tile\_height}$. The tile
set is the Cartesian product
\[
  T_k =
  \{\, yW+x \mid
      x \in [t^x_{\min},t^x_{\max}),\;
      y \in [t^y_{\min},t^y_{\max}) \,\}.
\]
If either radius component is non-positive, $T_k$ is empty.

\paragraph{AccuTile/SNUGBOX path.}
When both \texttt{conics} and \texttt{opacities} are present, the
kernel interprets
\[
  Q =
  \begin{bmatrix}
    A & B\\
    B & C
  \end{bmatrix}
\]
as the inverse projected covariance. It intersects tiles with the
opacity-thresholded ellipse
\[
  (p-\mu)^T Q (p-\mu) \le
  \min\!\left(3.33^2,\; 2\log(\alpha / \tau)\right),
  \qquad \tau = 1/255.
\]
The helper first computes a tight axis-aligned ellipse bounding box
(SNUGBOX), then scans tile columns or rows depending on the shorter
span. The fixture currently checks the AABB path against gsplat; the
AccuTile path is present in the vendored kernel and should get a
dedicated oracle fixture before relying on it in the projection
pipeline.

\subsection{Intersection-id encoding}

Each output intersection id is a signed 64-bit integer whose bit
layout matches gsplat:
\[
  \texttt{isect\_id}
  =
  (i \ll (32 + b_t))
  \;|\;
  (\texttt{tile\_id} \ll 32)
  \;|\;
  \operatorname{bits}_{32}(z),
\]
where $b_t = \lfloor \log_2(W H) \rfloor + 1$ and
$\operatorname{bits}_{32}(z)$ is the bit-level float32 representation
of depth, zero-extended to 64 bits. This ordering lets a later radix
sort group by image, then tile, then depth without carrying separate
side arrays.

\subsection{Launcher API \texorpdfstring{$\leftrightarrow$}{<->} math correspondence}

\begin{center}
\begin{tabular}{l l p{7.5cm}}
\hline
\textbf{C++ parameter} & \textbf{Math object} & \textbf{Meaning} \\
\hline
\texttt{packed} & --- & Selects dense $I\times N$ indexing or packed
$\mathrm{nnz}$ indexing. \\
\texttt{I}, \texttt{N}, \texttt{nnz} & $I,N,\mathrm{nnz}$ & Shape
parameters. \\
\texttt{image\_ids} & $i(p)$ & Required in packed mode; null in dense mode. \\
\texttt{means2d} & $\mu_k$ & Projected Gaussian centers. \\
\texttt{radii} & $r_k$ & Pixel-space AABB radii; non-positive component
culls the Gaussian. \\
\texttt{depths} & $z_k$ & Depth used in the low 32 bits of
\texttt{isect\_ids}. \\
\texttt{conics}, \texttt{opacities} & $Q_k,\alpha_k$ & Optional AccuTile
ellipse inputs. Both must be non-null to enable the ellipse path. \\
\texttt{cum\_tiles\_per\_gauss} & $C_k$ & Null for first pass; inclusive
prefix sum for second pass. \\
\texttt{tiles\_per\_gauss} & $c_k$ & First-pass output. \\
\texttt{isect\_ids}, \texttt{flatten\_ids} & --- & Second-pass outputs. \\
\texttt{stream} & --- & CUDA stream. \\
\hline
\end{tabular}
\end{center}

\subsection{Test correspondence}

The smoke tests in \texttt{IntersectTile\_tests.cu} cover:
\begin{enumerate}
  \item \emph{Dense AABB two-pass}: one Gaussian touches a $2\times2$
  tile block and one Gaussian is culled by zero radius. The test checks
  tile counts, encoded \texttt{isect\_ids}, and \texttt{flatten\_ids}.
  \item \emph{Packed AABB}: packed rows use explicit image ids in the
  high bits while keeping packed row indices as \texttt{flatten\_ids}.
\end{enumerate}
The oracle test in \texttt{IntersectTile\_oracle\_tests.cu} replays a
seeded gsplat-Python fixture captured with \texttt{isect\_tiles(...,
sort=False)} and checks both passes exactly. Sorting is intentionally
outside this launcher; a later CUB utility will own that boundary.

\section{Tile-offset encoding}
\label{sec:intersect-offset}

\paragraph{Source files.}
\begin{itemize}
  \item Vendored kernel (extracted from upstream gsplat):\\
  \texttt{Source/Kernels/Thirdparty/Gsplat/IntersectOffsetKernels.cuh}
  \item Raw-pointer launcher owned by heterosplat:\\
  \texttt{Source/Kernels/Heterosplat/IntersectOffset.\{h,cu\}}
  \item Tests and oracle fixtures:\\
  \texttt{Source/UnitTests/Kernels/Heterosplat/IntersectOffset\_tests.cu}\\
  \texttt{Source/UnitTests/Kernels/Heterosplat/IntersectOffset\_oracle\_tests.cu}\\
  \texttt{Source/UnitTests/Fixtures/IntersectOffset/}
\end{itemize}

\subsection{The forward map}

The tile-offset kernel converts a \emph{sorted} array of intersection
ids (output of \texttt{intersect\_tile} followed by radix sort) into a
per-image, per-tile offset array. The rasterizer later uses these
offsets to locate the Gaussians that contribute to each tile.

\paragraph{Input.}
A sorted array of $M$ intersection ids, each a signed 64-bit integer
whose upper 32~bits encode the composite key
\[
  k_j = i_j \cdot T + t_j,
\]
where $i_j$ is the image index, $t_j$ is the tile index within image
$i_j$, $T = W_t H_t$ is the total number of tiles per image, and
$W_t$, $H_t$ are the tile grid width and height. The lower 32~bits
carry depth and are ignored by this kernel. The array is sorted in
ascending order of $k_j$.

\paragraph{Output.}
An offset array $O$ of length $I \cdot T$, where $I$ is the number of
images. Each entry
\[
  O[i \cdot T + t] = \min \{ j : k_j \ge i \cdot T + t \}
\]
gives the index of the first intersection belonging to image~$i$,
tile~$t$ in the sorted intersection array. If no intersection exists
for that (image, tile) pair, the offset equals that of the next
non-empty slot, or $M$ for all trailing empty slots.

Equivalently, the half-open range of intersections for (image, tile)
pair $(i,t)$ is $[O[i T + t],\; O[i T + t + 1])$, where $O[I T]$
is implicitly~$M$.

\paragraph{Algorithm.}
Each CUDA thread handles one intersection~$j$. It decodes $k_j$ from
the upper 32 bits and, for the boundary cases:
\begin{enumerate}
  \item Thread 0 fills $O[0..k_0]$ with 0.
  \item Thread $M{-}1$ fills $O[k_{M-1}{+}1 .. IT{-}1]$ with $M$.
  \item All other threads compare $k_j$ to $k_{j-1}$; if they differ,
  fill $O[k_{j-1}{+}1 .. k_j]$ with $j$.
\end{enumerate}
This is a standard histogram-to-offset conversion (equivalent to an
exclusive scan of the histogram), executed in $O(M)$ total writes.
No backward pass is needed.

\subsection{Launcher API correspondence}

\begin{center}
\begin{tabular}{l l p{7.5cm}}
\hline
\textbf{C++ parameter} & \textbf{Math object} & \textbf{Meaning} \\
\hline
\texttt{n\_isects} & $M$ & Total number of sorted intersections. \\
\texttt{isect\_ids} & $\{k_j\}$ & Sorted intersection ids (upper 32
bits = composite key). \\
\texttt{I} & $I$ & Number of images. \\
\texttt{tile\_width} & $W_t$ & Tile grid width. \\
\texttt{tile\_height} & $H_t$ & Tile grid height. \\
\texttt{offsets} & $O$ & Output offset array, $[I \cdot T]$. \\
\texttt{stream} & --- & CUDA stream. \\
\hline
\end{tabular}
\end{center}

When $M = 0$ (no intersections), the launcher zeros the entire output
array via \texttt{cudaMemsetAsync} without launching the kernel.

\subsection{Test correspondence}

The smoke tests in \texttt{IntersectOffset\_tests.cu} cover:
\begin{enumerate}
  \item \emph{Single image, 2$\times$2 grid}: 3 intersections
  touching tiles 0, 1, 1. Verifies offsets $[0, 1, 3, 3]$.
  \item \emph{Multi-image}: 2 images, 2$\times$1 grid, verifying
  correct image-boundary encoding in the offset array.
  \item \emph{Zero intersections}: confirms the launcher zeros all
  offsets when the input is empty.
\end{enumerate}
The oracle test in \texttt{IntersectOffset\_oracle\_tests.cu} replays
a seeded gsplat-Python fixture captured with
\texttt{isect\_offset\_encode()} on sorted ids from
\texttt{isect\_tiles(..., sort=True)} and checks exact equality.

\section{EWA projection (fused 3D-to-2D)}
\label{sec:projection}

\paragraph{Source files.}
\begin{itemize}
  \item Vendored kernel templates (extracted from upstream gsplat):\\
  \texttt{Source/Kernels/Thirdparty/Gsplat/ProjectionEWA3DGSFusedKernels.cuh}
  (patched: \texttt{gpuAtomicAdd} $\to$ \texttt{atomicAdd}; ATen
  launchers and \texttt{GSPLAT\_BUILD\_3DGS} guard dropped).
  \item Raw-pointer launchers owned by heterosplat:\\
  \texttt{Source/Kernels/Heterosplat/ProjectionEWA3DGSFused.\{h,cu\}}.
  \item Tests and oracle fixtures:\\
  \texttt{Source/UnitTests/Kernels/Heterosplat/ProjectionEWA3DGSFused\_tests.cu}\\
  \texttt{Source/UnitTests/Kernels/Heterosplat/ProjectionEWA3DGSFused\_oracle\_tests.cu}\\
  \texttt{Source/UnitTests/Fixtures/ProjectionEWA3DGSFused/}
\end{itemize}

\subsection{The forward map}

The EWA projection kernel is the bridge between 3D Gaussians and the
2D tile rasterizer. For each Gaussian $n$ in batch $b$ viewed by
camera $c$, the kernel produces a projected 2D mean, an inverse 2D
covariance (the ``conic''), a depth, and pixel-space bounding radii.

\paragraph{Step 1: World to camera.}
Given a world-space mean $\mu_w \in \R^3$ and the camera extrinsics
$(R, t)$ extracted from a $4 \times 4$ row-major view matrix,
\begin{equation}
  \mu_c = R\, \mu_w + t.
  \label{eq:posw2c}
\end{equation}
If $\mu_c^z \notin [\texttt{near\_plane},\, \texttt{far\_plane}]$,
the Gaussian is culled (radii set to zero).

The 3D covariance transforms as
\begin{equation}
  \Sigma_c = R\, \Sigma_w\, R^T,
  \label{eq:covarw2c}
\end{equation}
where $\Sigma_w$ is either provided directly (triu-6 layout) or
computed from $(q, s)$ via the map in \S\ref{sec:qsc}.

\paragraph{Step 2: Perspective projection (pinhole).}
For a pinhole camera with intrinsics $f_x, f_y, c_x, c_y$, the
projected 2D mean is
\begin{equation}
  \mu_{2d} = \begin{pmatrix}
    f_x\, \mu_c^x / \mu_c^z + c_x \\
    f_y\, \mu_c^y / \mu_c^z + c_y
  \end{pmatrix}.
  \label{eq:persp-mean}
\end{equation}
The 2D covariance under the EWA (Elliptical Weighted Average)
approximation is
\begin{equation}
  \Sigma_{2d} = J\, \Sigma_c\, J^T,
  \label{eq:ewa}
\end{equation}
where $J \in \R^{2 \times 3}$ is the Jacobian of the pinhole map at
$\mu_c$:
\begin{equation}
  J = \frac{1}{\mu_c^z}
  \begin{pmatrix}
    f_x & 0 & -f_x\, \mu_c^x / \mu_c^z \\
    0 & f_y & -f_y\, \mu_c^y / \mu_c^z
  \end{pmatrix}.
  \label{eq:persp-jacobian}
\end{equation}
The kernel also supports orthographic and fisheye camera models via
a \texttt{camera\_model} enum, but the pinhole path is the one
exercised by all current tests.

\paragraph{Step 3: Anti-aliasing blur.}
To prevent singular conics when a Gaussian projects to sub-pixel
size, the kernel adds an isotropic blur:
\begin{equation}
  \Sigma_{2d}' = \Sigma_{2d} + \epsilon^2\, I_2,
  \label{eq:blur}
\end{equation}
where $\epsilon = \texttt{eps2d}$ (default $0.3$~px). The
compensation factor
\begin{equation}
  \kappa = \sqrt{\frac{\det \Sigma_{2d}}{\det \Sigma_{2d}'}},
  \label{eq:compensation}
\end{equation}
optionally written to \texttt{compensations}, can be applied to
opacity to preserve the integral of the Gaussian.

\paragraph{Step 4: Conic and radii.}
The conic (inverse 2D covariance) is
\begin{equation}
  Q = (\Sigma_{2d}')^{-1} =
  \begin{pmatrix} a & b \\ b & c \end{pmatrix},
  \label{eq:conic}
\end{equation}
stored as the upper triangle $(a, b, c)$. The tight pixel-space
bounding radii are
\begin{equation}
  r_x = \lceil \sigma_{\max} \sqrt{\Sigma_{2d,11}'} \rceil,
  \qquad
  r_y = \lceil \sigma_{\max} \sqrt{\Sigma_{2d,22}'} \rceil,
  \label{eq:radii}
\end{equation}
where $\sigma_{\max}$ is \texttt{GAUSSIAN\_EXTEND} ($3.33\sigma$) or
an opacity-aware tighter bound from~\cite{ye2024absgs}.

\subsection{Storage layout}

\begin{itemize}
  \item \texttt{means} --- $[B, N, 3]$ row-major world-space positions.
  \item \texttt{covars} --- $[B, N, 6]$ upper-triangle 3D covariance
  (optional; mutually exclusive with \texttt{quats}+\texttt{scales}).
  \item \texttt{quats}, \texttt{scales} --- $[B, N, 4]$ and $[B, N, 3]$
  (optional; mutually exclusive with \texttt{covars}).
  \item \texttt{viewmats} --- $[B, C, 4, 4]$ row-major extrinsics.
  \item \texttt{Ks} --- $[B, C, 3, 3]$ row-major intrinsics
  ($f_x = K_{00}$, $f_y = K_{11}$, $c_x = K_{02}$, $c_y = K_{12}$).
  \item \texttt{radii} --- $[B, C, N, 2]$ output, $(r_x, r_y)$ per
  Gaussian; $(0, 0)$ signals culled.
  \item \texttt{means2d} --- $[B, C, N, 2]$ projected centres.
  \item \texttt{depths} --- $[B, C, N]$ camera-space depth $\mu_c^z$.
  \item \texttt{conics} --- $[B, C, N, 3]$ upper triangle of $Q$.
\end{itemize}

\subsection{Launcher API \texorpdfstring{$\leftrightarrow$}{<->} math correspondence (forward)}

\begin{center}
\begin{tabular}{l l p{7.5cm}}
\hline
\textbf{C++ parameter} & \textbf{Math symbol} & \textbf{Meaning} \\
\hline
\texttt{B}, \texttt{C}, \texttt{N} & $B, C, N$ &
Batch size, camera count, Gaussian count. One CUDA thread per
$(b, c, n)$ triple. \\
\texttt{means} & $\{\mu_w\}$ & World-space means, $[B,N,3]$. \\
\texttt{covars} & $\{\Sigma_w\}$ & Optional triu-6 3D covariance. \\
\texttt{quats}, \texttt{scales} & $\{q\}, \{s\}$ & Alternative
covariance input via the $(q,s) \to \Sigma$ map (\S\ref{sec:qsc}). \\
\texttt{opacities} & $\{\alpha\}$ & Optional $[B,N]$; enables
opacity-aware radius tightening. \\
\texttt{viewmats} & $(R, t)$ & Camera extrinsics, $[B,C,4,4]$. \\
\texttt{Ks} & $f_x, f_y, c_x, c_y$ & Camera intrinsics, $[B,C,3,3]$. \\
\texttt{image\_width/height} & $W, H$ & Used for out-of-bounds culling. \\
\texttt{eps2d} & $\epsilon$ & Anti-aliasing blur radius~\eqref{eq:blur}. \\
\texttt{near\_plane}, \texttt{far\_plane} & --- & Depth culling range. \\
\texttt{radius\_clip} & --- & Minimum radius; Gaussians with both
$r_x, r_y \le$ this are culled. \\
\texttt{camera\_model} & --- & $0$=pinhole, $1$=ortho, $2$=fisheye. \\
\texttt{radii} & $r_x, r_y$ & Output bounding radii. \\
\texttt{means2d} & $\mu_{2d}$ & Output projected means. \\
\texttt{depths} & $\mu_c^z$ & Output depth. \\
\texttt{conics} & $Q$ & Output inverse 2D covariance (upper tri). \\
\texttt{compensations} & $\kappa$ & Optional blur compensation. \\
\texttt{stream} & --- & CUDA stream. \\
\hline
\end{tabular}
\end{center}

\subsection{The backward map (VJP)}
\label{sec:projection-bwd}

The backward kernel reverses the forward chain, threading gradients
$v_{\mu_{2d}}, v_z, v_Q$ (and optionally $v_\kappa$) back to the
inputs. The key VJP steps, in reverse order:

\paragraph{Conic inversion.}
For $Q = \Sigma^{-1}$, the matrix-inverse VJP gives
\begin{equation}
  v_{\Sigma_{2d}'} = -Q^T\, v_Q\, Q^T.
  \label{eq:inv-vjp}
\end{equation}
The kernel receives $Q$ from the forward pass (stored in
\texttt{conics}), so no matrix inversion is repeated.

\paragraph{EWA projection.}
The VJP of $\Sigma_{2d} = J \Sigma_c J^T$ and
$\mu_{2d} = \pi(\mu_c)$ produces $v_{\mu_c}$ and $v_{\Sigma_c}$.
For pinhole, the Jacobian~\eqref{eq:persp-jacobian} is
recomputed from the saved camera-space mean.

\paragraph{World-to-camera transform.}
From~\eqref{eq:posw2c} and~\eqref{eq:covarw2c}:
\begin{align}
  v_{\mu_w} &= R^T\, v_{\mu_c}, \\
  v_R &+= v_{\mu_c}\, \mu_w^T, \\
  v_{\Sigma_w} &= R^T\, v_{\Sigma_c}\, R, \\
  v_R &+= 2\, v_{\Sigma_c}\, R\, \Sigma_w.
  \label{eq:covar-vjp}
\end{align}

\paragraph{Quaternion+scale path.}
When the $(q,s)$ input was used, the kernel chains through the
\texttt{quat\_scale\_to\_covar\_vjp} helper (\S\ref{sec:qsc-bwd}),
producing $v_q$ and $v_s$ directly.

\paragraph{Warp-level reduction.}
Because multiple cameras may share the same Gaussian ($C > 1$), the
backward kernel uses warp-level labeled partitions
(\texttt{cg::labeled\_partition}) grouped by Gaussian~id, then a
single \texttt{atomicAdd} per warp group leader. This is why callers
must pre-zero all gradient output buffers.

\subsection{Launcher API \texorpdfstring{$\leftrightarrow$}{<->} math correspondence (backward)}

\begin{center}
\begin{tabular}{l l p{7.5cm}}
\hline
\textbf{C++ parameter} & \textbf{Math object} & \textbf{Meaning} \\
\hline
\texttt{B}, \texttt{C}, \texttt{N} & --- & Same shape parameters. \\
\texttt{means}, \texttt{covars}/\texttt{quats}/\texttt{scales} & --- &
Forward inputs, re-read to recompute intermediate values. \\
\texttt{viewmats}, \texttt{Ks} & $(R, t), K$ & Camera parameters. \\
\texttt{radii} & --- & Used to skip culled Gaussians. \\
\texttt{conics} & $Q$ & Forward output; used for the inverse VJP. \\
\texttt{compensations} & $\kappa$ & Forward output; needed if
\texttt{v\_compensations} is non-null. \\
\texttt{v\_means2d} & $\partial L / \partial \mu_{2d}$ &
Input gradient, $[B,C,N,2]$. \\
\texttt{v\_depths} & $\partial L / \partial z$ &
Input gradient, $[B,C,N]$. \\
\texttt{v\_conics} & $\partial L / \partial Q$ &
Input gradient, $[B,C,N,3]$. \\
\texttt{v\_compensations} & $\partial L / \partial \kappa$ &
Optional input gradient. \\
\texttt{v\_means} & $\partial L / \partial \mu_w$ &
Output gradient, $[B,N,3]$. Pre-zero; kernel accumulates. \\
\texttt{v\_covars} & $\partial L / \partial \Sigma_w$ &
Output gradient, $[B,N,6]$. Mutually exclusive with
\texttt{v\_quats}/\texttt{v\_scales}. \\
\texttt{v\_quats}, \texttt{v\_scales} & $\partial L / \partial q$,
$\partial L / \partial s$ &
Output gradients. Pre-zero; kernel accumulates. \\
\texttt{v\_viewmats} & $\partial L / \partial (R, t)$ &
Optional output gradient, $[B,C,4,4]$. Pre-zero; kernel accumulates. \\
\texttt{stream} & --- & CUDA stream. \\
\hline
\end{tabular}
\end{center}

\subsection{Numerical worth-noting}

\begin{itemize}
  \item The kernel recomputes $R(\hat q)$ and
  $\Sigma_w = R S^2 R^T$ in the backward pass rather than caching
  them from the forward. This is the same memory--compute tradeoff
  noted in \S\ref{sec:qsc-bwd-api}: the recompute is $O(N)$ and
  dominated by the rasterizer's per-tile work.
  \item The anti-aliasing blur~\eqref{eq:blur} adds $\epsilon^2 I_2$
  \emph{before} inversion, ensuring the conic is always
  well-conditioned. Without this, sub-pixel Gaussians produce
  singular $2 \times 2$ matrices.
  \item Gradient accumulation uses \texttt{atomicAdd} after warp-level
  reduction, so the number of atomic operations per Gaussian per
  output is $\lceil C / 32 \rceil$ in the worst case.
\end{itemize}

\subsection{Test correspondence}

The smoke tests in \texttt{ProjectionEWA3DGSFused\_tests.cu} cover:
\begin{enumerate}
  \item \emph{On-axis projection}: mean at $(0, 0, 5)$, identity
  view matrix, $f_x = f_y = 100$, $c_x = c_y = 64$. The projected
  mean should be $(64, 64)$ at depth $5$ with non-zero radii and
  finite conics.
  \item \emph{Behind-camera culling}: mean at $(0, 0, -1)$ is behind
  the camera; both radii must be zero.
  \item \emph{Backward smoke}: unit upstream gradients through a
  visible Gaussian; checks that $\partial L / \partial \mu_w$ is
  finite and non-zero.
\end{enumerate}
The oracle test in \texttt{ProjectionEWA3DGSFused\_oracle\_tests.cu}
replays a gsplat-Python fixture ($B=1$, $C=1$, $N=8$, pinhole) and
compares projected means, depths, and conics (masking culled
Gaussians) with $\text{atol} = 10^{-3}$, $\text{rtol} = 5 \times 10^{-3}$.

\section{Tile rasterizer (3DGS)}
\label{sec:rasterize}

\paragraph{Source files.}
\begin{itemize}
  \item Vendored kernel templates (extracted from upstream gsplat):\\
  \texttt{Source/Kernels/Thirdparty/Gsplat/RasterizeToPixels3DGSKernels.cuh}
  (patched: \texttt{gpuAtomicAdd} $\to$ \texttt{atomicAdd}; ATen
  launchers, \texttt{GSPLAT\_BUILD\_3DGS} guard, \texttt{Rasterization.h}
  include, and \texttt{GSPLAT\_FOR\_EACH} instantiations dropped).
  \item Raw-pointer launchers owned by heterosplat:\\
  \texttt{Source/Kernels/Heterosplat/RasterizeToPixels3DGS.\{h,cu\}}.
  \item Tests and oracle fixtures:\\
  \texttt{Source/UnitTests/Kernels/Heterosplat/RasterizeToPixels3DGS\_tests.cu}\\
  \texttt{Source/UnitTests/Kernels/Heterosplat/RasterizeToPixels3DGS\_oracle\_tests.cu}\\
  \texttt{Source/UnitTests/Fixtures/RasterizeToPixels3DGS/}
\end{itemize}

\subsection{The forward map}

The tile rasterizer is the final stage of the 3DGS rendering pipeline.
For each pixel $(i, j)$ in each image, it alpha-composites all
Gaussians that intersect the pixel's tile, in front-to-back depth
order. The inputs are the outputs of the earlier pipeline stages:
projected means and conics from the EWA projection
(\S\ref{sec:projection}), tile offsets from the intersection encoding
(\S\ref{sec:intersect-offset}), and per-Gaussian colours and
opacities.

\paragraph{Per-pixel compositing.}
For a pixel at position $p = (p_x, p_y)$ and the Gaussians
$g_1, \dots, g_K$ assigned to its tile (in depth-sorted order), the
compositing loop computes:
\begin{align}
  \sigma_k &= \tfrac{1}{2}\, \delta_k^T\, Q_k\, \delta_k,
  \qquad \delta_k = \mu_{2d,k} - p,
  \label{eq:rast-sigma}\\
  \alpha_k &= \min\!\bigl(\alpha_{\max},\; o_k\, e^{-\sigma_k}\bigr),
  \label{eq:rast-alpha}\\
  T_k &= \prod_{j=1}^{k-1} (1 - \alpha_j),
  \qquad T_1 = 1,
  \label{eq:rast-transmittance}\\
  C(p) &= \sum_{k=1}^{K} \alpha_k\, T_k\, c_k,
  \label{eq:rast-color}\\
  A(p) &= 1 - T_{K+1} = 1 - \prod_{k=1}^{K} (1 - \alpha_k),
  \label{eq:rast-alpha-out}
\end{align}
where $Q_k$ is the inverse 2D covariance (conic), $o_k$ is the base
opacity, $c_k \in \R^{\texttt{CDIM}}$ is the per-Gaussian colour, and
$\alpha_{\max}$ is a saturation cap (\texttt{MAX\_ALPHA} $= 0.999$).

The loop terminates early when the accumulated transmittance drops
below a threshold $\tau_T$ (\texttt{TRANSMITTANCE\_THRESHOLD}),
preventing negligible contributions from wasting compute.

\paragraph{Tile-cooperative loading.}
Each CUDA block processes one tile
(\texttt{tile\_size}~$\times$~\texttt{tile\_size} threads). Gaussians
are loaded from global memory into shared memory in batches of
\texttt{block\_size} by all threads cooperatively; each thread then
evaluates only its own pixel against the shared batch. This reduces
global memory traffic by a factor of $\texttt{tile\_size}^2$.

\paragraph{Background colour.}
If a per-image background $b \in \R^{\texttt{CDIM}}$ is provided, the
final colour includes the background blended through the remaining
transmittance:
\begin{equation}
  C_{\text{final}}(p) = C(p) + T_{K+1}\, b.
  \label{eq:rast-bg}
\end{equation}

\subsection{Storage layout}

\begin{itemize}
  \item \texttt{means2d} --- $[I, N, 2]$ or $[\text{nnz}, 2]$
  (packed mode). Interpreted as \texttt{vec2*}.
  \item \texttt{conics} --- $[I, N, 3]$ or $[\text{nnz}, 3]$. Upper
  triangle $(a, b, c)$ of the inverse 2D covariance; interpreted as
  \texttt{vec3*}.
  \item \texttt{colors} --- $[I, N, \texttt{CDIM}]$ or
  $[\text{nnz}, \texttt{CDIM}]$. Raw float array.
  \item \texttt{opacities} --- $[I, N]$ or $[\text{nnz}]$.
  \item \texttt{tile\_offsets} --- $[I, H_t, W_t]$ (int32). Start
  index into \texttt{flatten\_ids} for each tile.
  \item \texttt{flatten\_ids} --- $[M]$ (int32). Sorted Gaussian
  indices.
  \item \texttt{render\_colors} --- $[I, H, W, \texttt{CDIM}]$. Output.
  \item \texttt{render\_alphas} --- $[I, H, W]$. Output.
  \item \texttt{last\_ids} --- $[I, H, W]$ (int32). Index of last
  contributing Gaussian per pixel (used by backward).
\end{itemize}

\subsection{Launcher API \texorpdfstring{$\leftrightarrow$}{<->} math correspondence (forward)}

\begin{center}
\begin{tabular}{l l p{7.5cm}}
\hline
\textbf{C++ parameter} & \textbf{Math symbol} & \textbf{Meaning} \\
\hline
\texttt{I} & $I$ & Number of images. \\
\texttt{N} & $N$ & Gaussians per image (dense mode). \\
\texttt{n\_isects} & $M$ & Total sorted intersections. \\
\texttt{packed} & --- & Dense vs.\ packed indexing. \\
\texttt{means2d} & $\mu_{2d}$ & Projected centres. \\
\texttt{conics} & $Q$ & Inverse 2D covariance (conic). \\
\texttt{colors} & $c$ & Per-Gaussian colours. \\
\texttt{opacities} & $o$ & Per-Gaussian base opacities. \\
\texttt{backgrounds} & $b$ & Optional per-image background. \\
\texttt{masks} & --- & Optional per-tile skip mask. \\
\texttt{tile\_offsets} & --- & Per-tile start into sorted isects. \\
\texttt{flatten\_ids} & --- & Sorted Gaussian indices. \\
\texttt{render\_colors} & $C(p)$ & Output rendered colour. \\
\texttt{render\_alphas} & $A(p)$ & Output rendered alpha. \\
\texttt{last\_ids} & --- & Per-pixel last-Gaussian index (for bwd). \\
\texttt{stream} & --- & CUDA stream. \\
\hline
\end{tabular}
\end{center}

\subsection{The backward map (VJP)}

The backward kernel traverses Gaussians in \emph{back-to-front} order
(the reverse of the forward), recomputing $\alpha_k$ and $T_k$ on
the fly. The key gradient relations, derived from
\eqref{eq:rast-color}--\eqref{eq:rast-alpha-out}:

\paragraph{Colour gradient.}
\begin{equation}
  v_{c_k} = \alpha_k\, T_k\, v_{C(p)}.
  \label{eq:rast-bwd-color}
\end{equation}

\paragraph{Alpha gradient.}
The gradient $v_{\alpha_k} = \partial L / \partial \alpha_k$ mixes
contributions from the rendered colour, the rendered alpha, and
(optionally) the background:
\begin{equation}
  v_{\alpha_k} = \sum_{j} \Bigl(c_k^{(j)}\, T_k - B_k^{(j)}\, r_k\Bigr)\,
  v_{C(p)}^{(j)}
  + T_{\text{final}}\, r_k\, v_{A(p)},
  \label{eq:rast-bwd-alpha}
\end{equation}
where $r_k = 1 / (1 - \alpha_k)$ and $B_k$ is the accumulated colour
contribution from Gaussians \emph{behind} $k$.

\paragraph{Sigma and position gradients.}
From $\alpha_k = o_k\, e^{-\sigma_k}$:
\begin{align}
  v_{\sigma_k} &= -o_k\, e^{-\sigma_k}\, v_{\alpha_k},
  \label{eq:rast-bwd-sigma}\\
  v_{Q_k} &= \tfrac{1}{2}\, v_{\sigma_k}\,
  \begin{pmatrix} \delta_x^2 & 2\delta_x \delta_y & \delta_y^2 \end{pmatrix},
  \label{eq:rast-bwd-conic}\\
  v_{\mu_{2d,k}} &= v_{\sigma_k}\,
  \begin{pmatrix} Q_{11}\delta_x + Q_{12}\delta_y \\
  Q_{12}\delta_x + Q_{22}\delta_y \end{pmatrix}.
  \label{eq:rast-bwd-mean}
\end{align}

\paragraph{Warp-level reduction.}
Many pixels within a tile contribute gradients to the same Gaussian.
The kernel uses \texttt{cg::tiled\_partition<32>} warp reductions
followed by a single \texttt{atomicAdd} per warp leader. Callers
must pre-zero all gradient output buffers.

\subsection{Launcher API \texorpdfstring{$\leftrightarrow$}{<->} math correspondence (backward)}

\begin{center}
\begin{tabular}{l l p{7.5cm}}
\hline
\textbf{C++ parameter} & \textbf{Math object} & \textbf{Meaning} \\
\hline
\texttt{I}, \texttt{N}, \texttt{n\_isects} & --- & Same shape parameters. \\
\texttt{means2d}, \texttt{conics}, \texttt{colors}, \texttt{opacities} &
--- & Forward inputs, re-read for VJP. \\
\texttt{render\_alphas}, \texttt{last\_ids} & --- & Forward outputs
needed for back-to-front traversal. \\
\texttt{v\_render\_colors} & $v_{C(p)}$ & Input gradient. \\
\texttt{v\_render\_alphas} & $v_{A(p)}$ & Input gradient. \\
\texttt{v\_means2d} & $v_{\mu_{2d}}$ & Output gradient. Pre-zero;
kernel accumulates. \\
\texttt{v\_means2d\_abs} & $|v_{\mu_{2d}}|$ & Optional absolute-value
gradient (adaptive densification). \\
\texttt{v\_conics} & $v_Q$ & Output gradient. Pre-zero; accumulates. \\
\texttt{v\_colors} & $v_c$ & Output gradient. Pre-zero; accumulates. \\
\texttt{v\_opacities} & $v_o$ & Output gradient. Pre-zero; accumulates. \\
\texttt{stream} & --- & CUDA stream. \\
\hline
\end{tabular}
\end{center}

\subsection{Numerical worth-noting}

\begin{itemize}
  \item The colour channel count \texttt{CDIM} is a compile-time
  template parameter. heterosplat instantiates \texttt{CDIM=3} (RGB).
  Upstream gsplat supports a wider set via
  \texttt{GSPLAT\_FOR\_EACH}; extending heterosplat requires adding
  explicit instantiations in the launcher \texttt{.cu}.
  \item The forward uses \texttt{\_\_syncthreads\_count(done)} to
  exit early when the entire tile has saturated. This is critical
  for performance in scenes where many tiles reach
  $T \le \tau_T$ early.
  \item Shared memory per block is
  $\texttt{tile\_size}^2 \cdot (\texttt{sizeof(int32)} +
  2 \cdot \texttt{sizeof(vec3)})$ for the forward and additionally
  $\texttt{tile\_size}^2 \cdot \texttt{CDIM} \cdot
  \texttt{sizeof(float)}$ for the backward (colours are cached for
  the gradient computation). The launcher calls
  \texttt{cudaFuncSetAttribute} to request sufficient dynamic shared
  memory.
  \item The backward traverses in reverse order and recomputes
  transmittances from the saved \texttt{render\_alphas} (final $T$)
  and per-step $\alpha$ values. This avoids storing intermediate
  $T_k$ per pixel but costs one \texttt{fmaxf} guarded reciprocal
  per step.
\end{itemize}

\subsection{Test correspondence}

The smoke tests in \texttt{RasterizeToPixels3DGS\_tests.cu} cover:
\begin{enumerate}
  \item \emph{Single Gaussian}: one Gaussian at tile centre with
  isotropic conic and red colour. The centre pixel should have
  non-zero alpha and red colour contribution.
  \item \emph{Zero intersections}: with no Gaussians, all pixels
  should have zero alpha.
  \item \emph{Backward smoke}: unit upstream gradients through a
  visible Gaussian; checks that $\partial L / \partial c$ is finite
  and non-zero.
\end{enumerate}
The oracle test in \texttt{RasterizeToPixels3DGS\_oracle\_tests.cu}
replays a gsplat-Python fixture ($I=1$, $N=8$, $32 \times 32$
image) through the full pipeline and compares rendered colours and
alphas with $\text{atol} = 10^{-3}$, $\text{rtol} = 5 \times 10^{-3}$.

\section{Future sections}

This document is complete for Phase 0b. Future phases may add custom
kernels (e.g.\ \texttt{apply\_homogeneous\_transform\_kernel} in
Phase 2) which will be documented following the same template.

\end{document}
